\documentclass[10pt]{article}
\usepackage[margin=2.1cm]{geometry}
\usepackage[utf8]{inputenc}
\usepackage[T1]{fontenc}
\usepackage{amsmath,amssymb,booktabs}
\setlength{\parindent}{0pt}
\setlength{\parskip}{0.35em}
\pagestyle{empty}

\begin{document}
\small
\begin{center}
{\Large Conclusiones del Experimento 1}\
Trayectoria discreta admisible en $\Delta_3$
\end{center}

\textbf{Configuraci\'on.} Se fij\'o $X_0=(0.72,0.18,0.10)$, $P=(0.15,0.30,0.55)$, $\varepsilon=0.03$, $\alpha=0.35$, $R=0.60$ y horizonte $T=40$. El paso usado fue
\[
X_{t+1}=B_{L_t}(X_t K_{\varepsilon}^{\mathrm{sm}}),
\qquad
L_t=\exp\!\bigl(\alpha(\log P-\log Y_t)\bigr),
\quad Y_t=X_tK_{\varepsilon}^{\mathrm{sm}}.
\]

\textbf{Resultado num\'erico.} El potencial KL pas\'o de
\[
E_P(X_0)=0.8555667502
\]
a
\[
E_P(X_{40})=0.0027135384,
\]
con trayectoria mon\'otona para la pol\'itica elegida. El estado final fue
\[
X_{40}=(0.1744046671,\,0.3058505994,\,0.5197447335).
\]

\textbf{Lectura.} La suavizaci\'on markoviana no expuls\'o la trayectoria del interior del s\'implex y la correcci\'on bayesiana redujo de manera fuerte la divergencia respecto del objetivo. Num\'ericamente, este experimento muestra que el paso compuesto es operativo como mecanismo de control discreto y que el potencial KL funciona como observable natural de descenso.

\textbf{Conclusi\'on.} Bajo estos par\'ametros, el esquema discreto produce una trayectoria admisible estable, interior y fuertemente orientada hacia $P$. Es un buen candidato para la figura principal de la exposici\'on porque hace visible la composici\'on Markov + Bayes y su efecto geom\'etrico.

\end{document}